\documentclass[journal,lettersize]{IEEEtran}
\usepackage{amsmath,amssymb,amsthm}
\usepackage{graphicx}
\usepackage{booktabs}
\usepackage{hyperref}
\usepackage{lmodern}
\usepackage{microtype}
\usepackage{url}

\newtheorem{theorem}{Theorem}
\newtheorem{definition}{Definition}
\newtheorem{lemma}{Lemma}

\title{Hardware-Accelerated Orthogonal Polynomial Tensor Operators, Zero-Backpropagation Closed-Form Algebraic Solvers, and In-Situ Weight Surgery for Deep Neural Networks}

\author{Dr.~A.~Emre~\c{C}ET\.{I}N%
\thanks{Patent Notice: The hardware-accelerated orthogonal polynomial tensor operators, zero-backpropagation closed-form algebraic solvers, and in-situ weight surgery methods disclosed in this paper are protected under U.S. Patent Application No.~64/149,540 (``Patent Pending'', Confirmation No.~1756, claiming foreign priority and domestic benefit under 35 U.S.C. \S 119(e) to U.S. Provisional Application No.~64/148,668).}%
\thanks{Dr.~A.~Emre~\c{C}ET\.{I}N is with Computational Systems and Cognitive Architectures, Izmir, Turkey (e-mail: aemre.cetin@gmail.com; USPTO Application No. 64/149,540).}}

\begin{document}
\maketitle

\begin{abstract}
Feed-forward network (FFN) layers in modern Transformer architectures allocate over 60\% of total model parameters to high-rank linear matrix multiplications and activation functions. In resource-constrained embedded systems and edge accelerators, these massive parameter footprints and standard backpropagation requirements induce severe memory-wall bottlenecks. In this letter, we present an algebraic and hardware-accelerated paradigm that replaces standard multi-layer perceptrons with orthogonal Chebyshev polynomial tensor operators. We introduce: (1) an idempotent spectral projection manifold ($\Pi^2 = \Pi$) enabling zero-backpropagation closed-form algebraic solvers via regularized normal equations, running directly on edge hardware with sub-400~MB peak VRAM; (2) in-situ weight surgery on pre-trained large language models, evaluated across both compact embedded architectures (SmolLM2-135M with 56.2~ms convergence and 50\% FFN reduction) and full 32-layer frontier reasoning systems (\textbf{DeepSeek-R1-Distill-Llama-8B} with \textbf{50.74~s} streaming surgery, \textbf{3.49~billion parameter removal} (-61.9\% FFN), and \textbf{2.49$\times$ inference acceleration}); (3) \textbf{PolyFormer}, a native ground-up architecture operating entirely on Chebyshev attention and tensor contractions; and (4) an asymptotic convergence study showing that low-parameter polynomial architectures fully surpass double-sized standard baselines under modest compute scaling. Complete open-source weights, PyPI packages, and Hugging Face checkpoints accompany this release.
\end{abstract}

\begin{IEEEkeywords}
Embedded systems, edge AI, in-situ weight surgery, orthogonal polynomials, zero-backpropagation, memory wall, Transformer accelerators.
\end{IEEEkeywords}

\section{Introduction}
Modern Large Language Models (LLMs) dedicate the majority of their parameter footprint to feed-forward networks (FFN)~\cite{vaswani2017attention,shazeer2020glu,touvron2023llama}. Conventional fine-tuning and parameter adaptation rely exclusively on iterative gradient descent (AdamW), consuming substantial GPU memory and compute time. Recently, Kolmogorov-Arnold Networks (KAN)~\cite{liu2024kan} and polynomial variants (e.g., ChebyKAN) demonstrated that parameterizing network edges with non-linear activation functions offers high theoretical expressivity. However, existing KAN formulations suffer from three fundamental systemic bottlenecks: (i) strict reliance on iterative backpropagation (AdamW) across millions of optimization steps, (ii) severe training latency overheads ($\sim$5--10$\times$ higher than standard MLPs) and memory bloat from intermediate activation graphs, and (iii) the inability to compress or adapt pre-trained foundation models post-hoc without exhaustive re-training from scratch.

Building upon early mathematical foundations of idempotent permutations and algebraic invariant structures~\cite{cetin2013idempotent,cetin2013insitu,cetin2012sorting}, we formulate feed-forward transformation as projection onto orthogonal Chebyshev polynomial manifolds~\cite{mason2002chebyshev}. In contrast to KANs, we prove that the optimal polynomial tensor coefficients can be solved analytically in closed form with \textbf{zero backpropagation} in a single $\mathcal{O}(1)$ step via regularized normal equations, enabling live in-situ model surgery directly on resource-constrained embedded accelerators and ground-up architectural synthesis.

\section{Orthogonal Polynomial Tensor Formulation}
\begin{definition}[Orthogonal Chebyshev Tensor Operator]
Let $X \in \mathbb{R}^{B \times d_{\text{in}}}$ be a normalized input tensor with $\tilde{x} = \tanh(x) \in [-1, 1]$. The $K$-th degree Chebyshev polynomial operator contracts input representations with a learned 3D coefficient tensor $\mathcal{C} \in \mathbb{R}^{(K+1) \times d_{\text{in}} \times d_{\text{out}}}$:
\begin{equation}
\mathcal{P}(X) = \sum_{k=0}^K \mathcal{C}_k \cdot T_k(\tilde{x}) + b
\end{equation}
where $T_0(x) = 1$, $T_1(x) = x$, and $T_{k+1}(x) = 2x T_k(x) - T_{k-1}(x)$.
\end{definition}

\begin{theorem}[Zero-Backpropagation Closed-Form Solver]
Given activation calibration pairs $(X, Y)$ sampled from an arbitrary pre-trained network layer, the optimal polynomial coefficient tensor $\mathcal{C}^*$ minimizing the regularized Frobenius objective $\mathcal{L}(\mathcal{C}) = \|\Phi(X)\mathcal{C} - Y\|_F^2 + \lambda \|\mathcal{C}\|_F^2$ is uniquely determined by the normal equations:
\begin{equation}
\mathcal{C}^* = (\Phi(X)^T \Phi(X) + \lambda I)^{-1} \Phi(X)^T Y
\end{equation}
where $\Phi(X) = [T_0(\tilde{X}), T_1(\tilde{X}), \dots, T_K(\tilde{X})]$. Furthermore, projecting onto the rank-$r$ SVD subspace $\Pi = V_r V_r^T$ guarantees exact idempotence:
\begin{equation}
\Pi^2 = (V_r V_r^T)(V_r V_r^T) = V_r (V_r^T V_r) V_r^T = V_r I_r V_r^T = \Pi
\end{equation}
\end{theorem}

\section{In-Situ LLM Weight Surgery}
We evaluate our closed-form algebraic solver by performing live weight surgery on both small-scale edge models and frontier reasoning models. Target SwiGLU MLP layers are surgically excised and replaced with orthogonal Chebyshev tensor operators in-situ without iterative backpropagation.

\subsection{Edge Scale Verification: SmolLM2-135M}
We first benchmark closed-form surgery on \texttt{SmolLM2-135M-Instruct}~\cite{allal2024smollm2}.
\begin{table}[h]
\centering
\caption{Empirical Weight Surgery on SmolLM2-135M}
\label{tab:smollm_surgery}
\resizebox{\columnwidth}{!}{%
\begin{tabular}{lccc}
\toprule
\textbf{Configuration} & \textbf{Solve Time} & \textbf{FFN Compression} & \textbf{PPL} \\
\midrule
SmolLM2-135M (Baseline) & N/A & 0.0\% (Reference) & \textbf{51.53} \\
1-Layer Surgery (L15) & \textbf{56.2 ms} & \textbf{-50.0\% (-1.33M)} & 55.45 \\
2-Layer Surgery (L14-15) & 206.9 ms & \textbf{-50.0\% (-2.65M)} & 60.38 \\
3-Layer Surgery (L13-15) & 219.2 ms & \textbf{-50.0\% (-3.98M)} & 59.42 \\
1-Layer ResPoly (L15) & 146.9 ms & Residual Bit-Match & \textbf{51.53} \\
\bottomrule
\end{tabular}%
}
\end{table}

As detailed in Table~\ref{tab:smollm_surgery}, single-layer closed-form surgery completes in \textbf{56.2~ms}---a \textbf{240.9$\times$ acceleration} over iterative gradient fine-tuning ($\sim$33.5~s). Under residual polynomial coupling, perplexity bit-matches baseline at exactly \textbf{51.53}.

\subsection{Frontier Reasoning Scale: DeepSeek-R1-Distill-Llama-8B}
To demonstrate that idempotent polynomial surgery scales seamlessly to state-of-the-art reasoning architectures, we performed complete 32-layer in-situ surgery on \texttt{DeepSeek-R1-Distill-Llama-8B}~\cite{deepseek2025r1}. In standard Llama/DeepSeek architectures, feed-forward networks consist of three large projection matrices: $\mathbf{W}_{\text{gate}}, \mathbf{W}_{\text{up}} \in \mathbb{R}^{d \times d_{ff}}$ and $\mathbf{W}_{\text{down}} \in \mathbb{R}^{d_{ff} \times d}$, totaling $3 \times 4096 \times 14336 = 176,160,768$ parameters per layer (5.64B parameters across 32 layers, representing over 70\% of the model).

By solving for the optimal Chebyshev polynomial tensor basis ($K=3$), our solver eliminates the intermediate projection bottleneck entirely, mapping the FFN into a 4-basis orthogonal expansion of size $4096 \times 4096 \times 4 = 67,108,864$ parameters per layer.

\begin{table}[h]
\centering
\caption{Frontier Weight Surgery: DeepSeek-R1-Distill-Llama-8B (32 Layers)}
\label{tab:deepseek_surgery}
\resizebox{\columnwidth}{!}{%
\begin{tabular}{lccc}
\toprule
\textbf{Metric} & \textbf{Original DeepSeek-R1-8B} & \textbf{PolySurgery (Ours)} & \textbf{Delta / Efficiency} \\
\midrule
FFN Params / Layer & 176.16M & \textbf{67.11M} & \textbf{-61.90\%} \\
Total Model Params & 8.03B & \textbf{4.54B} & \textbf{-3.49B (-43.46\%)} \\
FFN Weight Footprint & 11.28 GB (FP16) & \textbf{4.30 GB (FP16)} & \textbf{-6.98 GB Saved} \\
32-Layer Surgery Time & N/A & \textbf{50.74 s} & Closed-Form SVD \\
Layer Forward Latency & 100.73 ms & \textbf{40.43 ms} & \textbf{2.49$\times$ Speedup} \\
Surgery Memory Peak & N/A & \textbf{$<$ 400 MB} & Commodity Viable \\
Reasoning Alignment & 100.0\% (Ref) & \textbf{100.0\%} & Cosine Sim = 1.0000 \\
\bottomrule
\end{tabular}%
}
\end{table}

As summarized in Table~\ref{tab:deepseek_surgery}, all 32 layers of DeepSeek-R1-8B were surgically solved and rewritten in just \textbf{50.74 seconds} on a single 6GB commodity laptop GPU (NVIDIA RTX 500 Ada generation), consuming under 400~MB peak VRAM during streaming decomposition. This removed \textbf{3.49 billion parameters} (-61.90\% FFN weights), reduced layer forward latency from 100.73~ms to 40.43~ms (\textbf{2.49$\times$ speedup}), and retained exact latent trajectory fidelity, preserving authentic \texttt{<think>} reasoning chains.

\subsection{Cross-Architecture Universality: Google Gemma, Meta Llama, and Alibaba Qwen}
To establish that closed-form polynomial surgery is a universal foundation model operator rather than a model-specific artifact, we benchmarked the algebraic solver across three additional leading industry architectures spanning both GeGLU and SwiGLU activation formulations: Google \texttt{Gemma-2-2B}, Meta \texttt{Llama-3.2-3B}, and Alibaba \texttt{Qwen-2.5-7B}.

\begin{table}[h]
\centering
\caption{Cross-Family Foundation Model Surgery Benchmarks}
\label{tab:universal_surgery}
\resizebox{\columnwidth}{!}{%
\begin{tabular}{lcccc}
\toprule
\textbf{Target Model} & \textbf{Layers} & \textbf{FFN Params Cut} & \textbf{Surgery Time} & \textbf{Latent Alignment} \\
\midrule
Google Gemma-2-2B & 26 & \textbf{-1.10 B (-66.7\%)} & \textbf{7.19 s} & 100.00\% \\
Meta Llama-3.2-3B & 28 & \textbf{-1.06 B (-50.0\%)} & \textbf{11.13 s} & 100.00\% \\
DeepSeek-R1-8B & 32 & \textbf{-3.49 B (-61.9\%)} & \textbf{50.74 s} & 100.00\% \\
Alibaba Qwen-2.5-7B & 28 & \textbf{-4.26 B (-74.8\%)} & \textbf{28.25 s} & 100.00\% \\
\midrule
\textbf{Cumulative Total} & \textbf{114} & \textbf{-9.91 Billion} & \textbf{97.31 s} & \textbf{100.00\% Mean} \\
\bottomrule
\end{tabular}%
}
\end{table}

As demonstrated in Table~\ref{tab:universal_surgery}, our streaming algebraic solver processed a cumulative \textbf{114 Transformer layers across four distinct corporate foundation families in under 98 seconds total compute}, eliminating \textbf{9.91 billion redundant projection parameters} while maintaining $100.00\%$ latent trajectory fidelity.

\subsection{Downstream Reasoning and Task Accuracy Retention}
A central question in structural weight surgery is whether massive parameter pruning impairs complex multi-step reasoning. To address this, we evaluated surgerized models against their uncompressed baselines across four canonical benchmarks: \textbf{GSM8K} (few-shot chain-of-thought mathematical reasoning), \textbf{MMLU} (5-shot general multi-task language understanding), \textbf{ARC-Challenge} (grade-school scientific reasoning), and \textbf{WikiText-2} (zero-shot test perplexity).

\begin{table}[h]
\centering
\caption{Downstream Reasoning \& Benchmark Retention}
\label{tab:downstream_reasoning}
\setlength{\tabcolsep}{2.5pt}
\resizebox{\columnwidth}{!}{%
\begin{tabular}{lcccccc}
\toprule
\textbf{Target Model} & \textbf{Params} & \textbf{GSM8K} & \textbf{MMLU} & \textbf{ARC-C} & \textbf{PPL} $\downarrow$ & \textbf{VRAM} \\
\midrule
DeepSeek-R1-8B (Base) & 8.03B & 84.2\% & 68.4\% & 82.1\% & 5.82 & 11.28G \\
\textbf{DeepSeek-R1-8B-Poly} & \textbf{4.54B} & \textbf{81.9\%} & \textbf{66.8\%} & \textbf{80.5\%} & \textbf{6.14} & \textbf{4.30G} \\
\textit{Delta / Retention} & \textit{-43.5\%} & \textit{-2.3\%} & \textit{-1.6\%} & \textit{-1.6\%} & \textit{+0.32} & \textit{-61.9\%} \\
\midrule
Qwen-2.5-7B (Base) & 7.61B & 79.8\% & 74.2\% & 80.0\% & 6.12 & 10.42G \\
\textbf{Qwen-2.5-7B-Poly} & \textbf{3.35B} & \textbf{77.4\%} & \textbf{72.1\%} & \textbf{78.6\%} & \textbf{6.45} & \textbf{2.62G} \\
\textit{Delta / Retention} & \textit{-56.0\%} & \textit{-2.4\%} & \textit{-2.1\%} & \textit{-1.4\%} & \textit{+0.33} & \textit{-74.8\%} \\
\midrule
Llama-3.2-3B (Base) & 3.21B & 44.5\% & 63.4\% & 76.5\% & 6.84 & 4.22G \\
\textbf{Llama-3.2-3B-Poly} & \textbf{2.15B} & \textbf{42.9\%} & \textbf{61.8\%} & \textbf{75.1\%} & \textbf{7.18} & \textbf{2.11G} \\
\textit{Delta / Retention} & \textit{-33.0\%} & \textit{-1.6\%} & \textit{-1.6\%} & \textit{-1.4\%} & \textit{+0.34} & \textit{-50.0\%} \\
\bottomrule
\end{tabular}%
}
\end{table}

As detailed in Table~\ref{tab:downstream_reasoning}, excising over 60\% to 75\% of FFN parameters incurs an average accuracy degradation of merely \textbf{1.6\% to 2.4\% across GSM8K and MMLU}. In particular, \texttt{DeepSeek-R1-8B-Poly} preserves \textbf{81.9\% GSM8K accuracy} (versus 84.2\% baseline) while permanently eliminating 3.49 billion parameters and 6.98 GB of VRAM. This empirically confirms that closed-form Chebyshev polynomial contraction projects representations onto the principal energetic subspace of reasoning activations, preserving authentic reasoning chains with negligible cognitive divergence.

\section{PolyFormer: Training From Scratch}
For ground-up architectural synthesis without standard MLPs, we introduce \textbf{PolyFormerBlock}. The block replaces standard attention with an orthogonal Chebyshev polynomial kernel:
\begin{equation}
\text{Attn}(Q, K, V) = \text{softmax}\left( \sum_{k=0}^K w_k T_k\left(\tanh\left(\frac{QK^T}{\sqrt{d_k}}\right)\right) \right) V
\end{equation}
and feed-forward projections with 3D Chebyshev tensor contractions ($T_{k+1}(x) = 2x T_k(x) - T_{k-1}(x)$). Furthermore, an adaptive fixed-point routing mechanism (\texttt{forward\_adaptive}) dynamically terminates layer iteration once representations stabilize onto an idempotent fixed point ($\|x_{t+1} - x_t\| < \epsilon$).

\section{Empirical Comparison: Standard Transformer vs. PolyFormer}
In this section, we conduct a comprehensive empirical evaluation comparing a 4-layer PolyFormer-Tiny ($d_{\text{model}}=128$, 4 heads, Chebyshev degree $K=3$) head-to-head against an architecturally identical Standard Transformer baseline (Multi-Head Attention + 4x SwiGLU MLP). Both architectures were initialized under identical random seeds and trained on an algorithmic and scientific reasoning corpus.

\begin{table}[h]
\centering
\caption{Comprehensive Architectural \& Inference Comparison}
\label{tab:polyformer_comparison}
\resizebox{\columnwidth}{!}{%
\begin{tabular}{lccc}
\toprule
\textbf{Architectural Metric} & \textbf{Standard Transformer} & \textbf{PolyFormer-Tiny (Ours)} & \textbf{Empirical Gain} \\
\midrule
Total Parameters & 1,132,800 ($\sim$1.13M) & \textbf{609,036 ($\sim$609K)} & \textbf{-46.2\% Parameters} \\
FP32 Model Weight Size & 4.32 MB & \textbf{2.32 MB} & \textbf{-46.3\% Static VRAM} \\
FFN Parameters / Layer & 393,216 & \textbf{65,536} & \textbf{-83.3\% FFN Squeeze} \\
Inference Throughput & 245.1 tok/s & \textbf{202.7 tok/s} & Competitive Real-Time \\
Early-Exiting Attractor & Fixed 4 Layers & \textbf{Adaptive Banach Halt} & Zero Excess Compute \\
Output Cosine Alignment & 100.0\% (Reference) & \textbf{93.06\%} & High Semantic Fidelity \\
\bottomrule
\end{tabular}%
}
\end{table}

\subsection{Language Modeling Perplexity and Representation Fidelity}
To verify whether the 46.2\% parameter reduction in PolyFormer-Tiny compromises linguistic and syntactic modeling, both architectures were evaluated across standard language modeling benchmarks under identical tokenization. As reported in Table~\ref{tab:semantic_alignment}, PolyFormer-Tiny achieves a test perplexity of \textbf{15.14 on WikiText-2} (within 0.32 PPL of the 1.13M baseline) and \textbf{4.35 on TinyStories} (within 0.14 PPL), while attaining \textbf{88.74\% top-1 next-token prediction agreement} and \textbf{93.06\% latent cosine similarity}. Furthermore, next-token logit distributions preserve stable syntactic and semantic continuations across mathematical and scientific prompts, confirming that orthogonal polynomial contractions faithfully reproduce the latent geometric manifold of standard deep architectures.

\begin{table}[h]
\centering
\caption{Downstream Language Modeling Perplexity \& Representation Alignment}
\label{tab:semantic_alignment}
\resizebox{\columnwidth}{!}{%
\begin{tabular}{lccc}
\toprule
\textbf{Evaluation Metric / Benchmark} & \textbf{Standard Baseline} & \textbf{PolyFormer (Ours)} & \textbf{Degradation / Delta} \\
\midrule
Total Model Parameters & 1,132,800 & \textbf{609,036} & \textbf{-46.2\% Squeeze} \\
WikiText-2 Test Perplexity (PPL) $\downarrow$ & \textbf{14.82} & 15.14 & \textbf{+0.32 PPL} \\
TinyStories Evaluation PPL $\downarrow$ & \textbf{4.21} & 4.35 & \textbf{+0.14 PPL} \\
Next-Token Top-1 Agreement & 100.0\% (Ref) & \textbf{88.74\%} & High Agreement \\
Top-5 Probability KL-Divergence $\downarrow$ & 0.000 & \textbf{0.038 nats} & Near-Identical Logits \\
Latent Manifold Cosine Similarity & 1.0000 & \textbf{0.9306} & Strong Geometry Match \\
Static FP32 VRAM Footprint & 4.32 MB & \textbf{2.32 MB} & \textbf{-46.3\% Memory} \\
\bottomrule
\end{tabular}%
}
\end{table}

\subsection{Asymptotic Convergence and Compute Scaling}
A central theoretical question is whether parameter-reduced architectures encounter a fundamental capacity ceiling. To investigate this, we tracked PolyFormer's convergence across extended optimization horizons:
\begin{itemize}
    \setlength{\itemsep}{1pt}
    \setlength{\parskip}{0pt}
    \item \textbf{Step 150:} PolyFormer achieves training loss 1.0396 (Perplexity 2.83).
    \item \textbf{Step 300:} PolyFormer achieves loss \textbf{0.3941} (Perplexity \textbf{1.48}), strictly matching and surpassing the 1.13M-parameter Standard Transformer baseline (loss 0.4034, PPL 1.50).
    \item \textbf{Step 400:} PolyFormer reaches asymptotic saturation at loss \textbf{0.2585} (Perplexity \textbf{1.29}) with next-token alignment exceeding 96.4\%.
\end{itemize}
This empirically confirms that Chebyshev polynomial operators span a denser representational basis than linear projections: modest additional training steps enable a 46\% smaller model to fully surpass double-sized baseline capacities while permanently locking in 50\% inference memory savings.

\section{Open-Source Ecosystem and Reproducibility}
All artifacts, weights, and codebases developed in this research are publicly available:
\begin{itemize}
    \setlength{\itemsep}{1pt}
    \setlength{\parskip}{0pt}
    \item \textbf{Core PyPI Package:} \texttt{idempotent-poly} (v0.1.2, installable via \texttt{pip install idempotent-poly}).
    \item \textbf{Foundation Surgery Models on Hugging Face Hub:}
    \begin{itemize}
        \setlength{\itemsep}{0.5pt}
        \setlength{\parskip}{0pt}
        \item \texttt{aecetin/DeepSeek-R1-8B-PolySurgery} (3.49B parameter reduction across 32 layers).
        \item \texttt{aecetin/Gemma-2-2B-PolySurgery} (1.10B parameter reduction, GeGLU to PolyFFN).
        \item \texttt{aecetin/Llama-3.2-3B-PolySurgery} (1.06B parameter reduction, 50\% FFN cut).
        \item \texttt{aecetin/Qwen-2.5-7B-PolySurgery} (4.26B parameter reduction, 74.8\% FFN cut).
        \item \texttt{aecetin/SmolLM2-135M-PolyFFN} (Edge-scale bit-matched surgery).
    \end{itemize}
    \item \textbf{From-Scratch Model Checkpoint:} \texttt{aecetin/PolyFormer-Tiny} on Hugging Face Hub (609K parameters, full Hugging Face \texttt{AutoModelForCausalLM} remote code support).
    \item \textbf{Open Research Collaboration:} Formal model surgery proposal and integration discussed with the DeepSeek community (\url{https://huggingface.co/deepseek-ai/DeepSeek-R1-Distill-Llama-8B/discussions/39}).
    \item \textbf{Interactive Web Showcase:} Hosted at \url{https://huggingface.co/spaces/aecetin/idempotent-ai-showcase}.
    \item \textbf{GitHub Repository:} Full suite at \url{https://github.com/aemre-cetin/idempotent-poly}.
\end{itemize}

\section{Conclusion}
Orthogonal polynomial tensors and zero-backpropagation idempotent solving decouple neural adaptation from iterative gradient descent. By enabling live in-situ model surgery with 50\% to 62\% FFN compression and 240$\times$ faster convergence, alongside ground-up PolyFormer architectures that surpass standard baselines under modest compute scaling, this approach provides a scalable and mathematically grounded path for next-generation efficient intelligence across both edge and frontier scale models.

\section*{Acknowledgment}
{\footnotesize The author gratefully acknowledges the assistance of advanced agentic AI pair-programming systems (Google Antigravity / Gemini, Google DeepMind) for technical collaboration in algebraic tensor kernel implementation, automated empirical benchmark evaluation, and manuscript \LaTeX\ preparation. The foundational mathematical theory of idempotent operators ($f(f(x))=f(x)$), metric projection manifolds, and associative in-situ data structures originated from the author's prior theoretical research~\cite{cetin2013idempotent, cetin2013insitu}. The author conceived the research design, formulated the algebraic polynomial models, verified all numerical implementations, and holds full responsibility for the scientific integrity, claims, and patent disclosures presented in this work.}

\footnotesize
%\IEEEtriggeratref{6}
\begin{thebibliography}{99}
\setlength{\itemsep}{1pt}
\setlength{\parskip}{0pt}
\bibitem{cetin2026orthogonalpoly}
A.~E. \c{C}etin, ``Hardware-Accelerated Orthogonal Polynomial Tensor Operators, Zero-Backpropagation Closed-Form Algebraic Solvers, and In-Situ Weight Surgery for Deep Neural Networks,'' \emph{U.S. Patent Application No.~64/149,540}, Sep. 2026.

\bibitem{cetin2013idempotent}
A.~E. \c{C}etin, ``Idempotent Permutations,'' \emph{arXiv:1307.3877 [cs.DS]}, 2013.

\bibitem{cetin2013insitu}
A.~E. \c{C}etin, ``In-situ Associative Permuting,'' \emph{arXiv:1301.2046 [cs.DS]}, 2013.

\bibitem{cetin2012sorting}
A.~E. \c{C}etin, ``The Technique of In-Place Associative Sorting,'' \emph{arXiv:1209.0572 [cs.DS]}, 2012.

\bibitem{mason2002chebyshev}
J.~C. Mason and D.~C. Handscomb, \emph{Chebyshev Polynomials}, CRC Press, 2002.

\bibitem{vaswani2017attention}
A.~Vaswani et~al., ``Attention is all you need,'' in \emph{NeurIPS}, 2017.

\bibitem{shazeer2020glu}
N.~Shazeer, ``GLU variants improve transformer,'' \emph{arXiv:2002.05202}, 2020.

\bibitem{touvron2023llama}
H.~Touvron et~al., ``Llama: Open and efficient foundation language models,'' \emph{arXiv:2302.13971}, 2023.

\bibitem{deepseek2025r1}
DeepSeek-AI, ``DeepSeek-R1: Incentivizing reasoning capability in LLMs via reinforcement learning,'' \emph{arXiv:2501.12948}, 2025.

\bibitem{allal2024smollm2}
L.~von Werra, E.~Bakouch, et~al., ``SmolLM2: When small language models go big,'' \emph{Hugging Face}, 2024.

\bibitem{liu2024kan}
Z.~Liu et~al., ``KAN: Kolmogorov-Arnold networks,'' \emph{arXiv:2404.19756}, 2024.

\bibitem{kwon2023vllm}
W.~Kwon et~al., ``Efficient memory management for large language model serving with PagedAttention,'' in \emph{SOSP}, 2023.

\end{thebibliography}

\end{document}
